\subsection{Problem Statement}

We begin with the simplest non-trivial formulation of our problem, inspired by \ref{ex:motivating-problem}.

Consider a graph $G = (V,E)$ and subgraphs $G_1, \ldots, G_k \subset G$. Let $G_i = (V_i, E_i)$ for these subgraphs. For now let us suppose $G$ is unweighted and undirected. The subgraphs need not cover $G$ and can overlap. Each player $i$ knows $G$ and $G_i$, but not the other player's graphs. The goal is, given some starting vertex $v \in V$, and some end vertex $w \in V$ to:
\begin{enumerate}
	\item  Identify if it is possible to construct a path $v, v_1, \dots, v_l, w$ such that every edge on the path is in $\bigcup_{i=0}^{k} E_i$. \label{prob:traversal-exists-path}
	\item  If such a path exists, find the minimum number of edges needed to get from $v$ to $w$, while using only edges belonging to $\bigcup_{i=0}^{k} E_i$. \label{prob:traversal-min-path}
\end{enumerate}
If some $e = (u,v) \in G_i$, we must have $u,v \in V(G_i)$. But the converse is not always true.

\begin{figure}[H]
\centering
\begin{tikzpicture}
  \draw (0,0) ellipse [x radius=6, y radius=2.5];
  \node at (0,2.1) {$G$};

  \draw plot [smooth cycle, tension=0.6] coordinates {(-2.2,-0.5) (-1.2,0.4)
	  (1.2,0.4) (1.8,-0.8) (0.2,-1.6) (-1.6,-1.7) };
  \node at (-1.1,-1.3) {$G_1$};

  \draw plot [smooth cycle, tension=1.1] coordinates {(-0.6,1.2) (0.8,1.6)
	  (2.1,0.6) (1.1,0) (-0.4,0.3)};
  \node at (0.7,1.2) {$G_2$};

  \draw plot [smooth cycle, tension=0.8] coordinates {(2.3,0.9) (4.6,1.2)
	  (4.4,-0.2) (3.2,-0.9) (1.0,-0.2)};
  \node at (4.2,0.8) {$G_3$};

  \node[point] (a1) at (-1.5,0.1) {};
  \node[point] (a2) at (-0.5,-0.6) {};
  \node[point] (a3) at (0.8,-0.4) {};
  \draw[dotted] (a1) -- (a2) -- (a3);

  \node[point] (b1) at (-0.2,0.9) {};
  \node[point] (b2) at (0.8,0.7) {};
  \node[point] (b3) at (1.4,0.2) {};
  \draw[dotted] (b1) -- (b2) -- (b3);

  \node[point] (c1) at (2.6,0.5) {};
  \node[point] (c2) at (3.8,0.6) {};
  \node[point] (c3) at (3.4,-0.3) {};
  \draw[dotted] (c1) -- (c2) -- (c3);

  \node[point] (d1) at (-5,0) {};
  \node[point] (d2) at (-3,-0.7) {};
  \node[point] (d3) at (-2.4,1) {};
  \node[point] (d4) at (-4,0.4) {};
  \draw[dotted] (d1) -- (d2) -- (d4);
  \draw[dotted] (d2) -- (d3) -- (b1);

  \node[below left=0.1pt , outer sep=2pt] at (0.8,-0.4) {$v$};
  \node[above right=1pt , outer sep=2pt] at (2.6,0.5) {$w$};

  \draw[thick] (c1) -- (b3) -- (a3);
\end{tikzpicture}
\begin{quote}
    \caption{Example graph where $TRAV(V,W)=1$}
\end{quote}
\label{fig:traversal-valid-path}
\end{figure}

Sometimes, condition \ref{prob:traversal-exists-path} fails, as in Figure \ref{fig:traversal-no-path}. In such a case, it is sufficient for one party to be convinced that no path exists.

\begin{figure}[H]
\centering
\begin{tikzpicture}
  \draw (0,0) ellipse [x radius=6, y radius=2.5];
  \node at (0,2.1) {$G$};

  \draw plot [smooth cycle, tension=0.9] coordinates {(-4.2,0.5) (-3.2,1.4) (-1.8,0.8) (-2.2,-0.6) (-3.6,-0.7)};
  \node at (-3.1,1) {$G_A$};

  \draw plot [smooth cycle, tension=0.8] coordinates {(2.3,0.9) (4.6,1.2) (4.4,-0.2) (3.2,-0.9) (2.0,-0.2)};
  \node at (4.2,0.8) {$G_B$};

  \node[point] (a1) at (-3.6,0.2) {};
  \node[point] (a2) at (-2.6,0.6) {};
  \node[point] (a3) at (-2.8,-0.4) {};
  \draw[dotted] (a1) -- (a2) -- (a3);

  \node[point] (b1) at (-0.2,0.9) {};
  \node[point] (b2) at (0.8,0.7) {};
  \node[point] (b3) at (1.4,0.2) {};
  \draw[dotted] (b1) -- (b2) -- (b3);
  \draw[dotted] (a2) -- (b1);
  \draw[dotted] (b3) -- (c1);

  \node[point] (c1) at (2.6,0.5) {};
  \node[point] (c2) at (3.8,0.6) {};
  \node[point] (c3) at (3.4,-0.3) {};
  \draw[dotted] (c1) -- (c2) -- (c3);
\end{tikzpicture}
\caption{Example graph where $\TRAV(x,y) = 0$.}
\label{fig:traversal-no-path}
\end{figure}

\subsection{Two-party lower bounds}

We begin by considering the case with only two players and thus two subgraphs. Let these players be named Alice and Bob, with subgraphs $G_A$ and $G_B$ respectively. Denote this problem by $\TRAV_2$, i.e. the traversal problem for 2 players, where $|V(G)| = n$ and $|E(G)| = O(n^2)$.

\begin{claim}
\label{claim:equality-trav-bound}
Recall that equality $\EQ$ has inputs $x,y \in \{0,1\}^n$, and $\EQ(x,y) = 1$ iff $x = y$. Then $D(\TRAV_2) \geq D(\EQ)$ and so on.
\end{claim}

\begin{proof}
\textbf{Equality gadget.}
We construct a gadget that enforces equality at a single bit position. The gadget is shown in Figure~\ref{fig:equality_gadget}. The gadget consists of two parallel branches.

\begin{enumerate}
    \item The top branch corresponds to the assignment $x_i = y_i = 1$, 
    \item The bottom branch corresponds to the assignment $x_i = y_i = 0$. 
\end{enumerate}
 
\begin{figure}[H]
\centering
\begin{tikzpicture}[scale=1.2, every node/.style={circle, fill=black, inner sep=1.5pt}]

\node (L) at (-2,0) {};
\node (T1) at (-0.8,1.5) {};
\node (T2) at (0.8,1.5) {};
\node (B1) at (-0.8,-1.5) {};
\node (B2) at (0.8,-1.5) {};
\node (R) at (2,0) {};

\draw (L) -- (T1);
\draw (T1) -- (T2);
\draw (T2) -- (R);

\draw (L) -- (B1);
\draw (B1) -- (B2);
\draw (B2) -- (R);

\node[draw=none, fill=none, above=6pt] at (T1) {$x_i = 1$};
\node[draw=none, fill=none, above=6pt] at (T2) {$y_i = 1$};

\node[draw=none, fill=none, below=6pt] at (B1) {$x_i = 0$};
\node[draw=none, fill=none, below=6pt] at (B2) {$y_i = 0$};

\node[draw=none, fill=none, left=6pt] at (L) {$L_i$};
\node[draw=none, fill=none, right=6pt] at (R) {$R_i$};

\end{tikzpicture}
\caption{Gadget for proving equality. A path only exists if $x_i = y_i$.}
\label{fig:equality_gadget}
\end{figure}

The idea is that Alice enables exactly one entry node of the gadget depending on $x_i$, and Bob enables exactly one exit node depending on $y_i$. A path exists through the gadget if and only if Alice and Bob select nodes on the same branch, which holds exactly when $x_i = y_i$.

\textbf{Reduction to our graph problem.}
Construct a graph of $n$ equality gadgets in a chain, as in Figure \ref{fig:equality-gadget-chain}. Construct subgraph $G_A$ as follows: Include all vertices at the start and end of each equality gadget. Also, for every $i \in [n]$, give the vertex labelled $x_i = 1$ if this is true, else $x_i = 0$. Then, add all edges between these vertices, as well as all edges between $x_i = z$ and $y_i = z$ for $z \in \{0,1\}$. Do similar for $G_B$.

In the resulting graph, there exists a path from $v$ to $w$ if and only if all gadgets are traversable, which occurs if and only if $x_i = y_i$ for all $i$, i.e., if and only if $x = y$.

\begin{figure}[H]
\centering
\begin{tikzpicture}[
hex/.style={
draw,
regular polygon,
regular polygon sides=6,
minimum size=1.4cm,
inner sep=0pt
}
]

% Hexagons placed manually so edges touch
\node[hex] (h1) at (0,0) {};
\node[hex] (h2) at (1.4,0) {};
\node[hex] (h3) at (2.8,0) {};
\node[hex] (h4) at (4.2,0) {};
\node[hex] (h5) at (5.6,0) {};

% Optional labels
\node at (h1) {\scriptsize $i=1$};
\node at (h2) {\scriptsize $i=2$};
\node at (h3) {\scriptsize $i=3$};
\node at (h4) {\scriptsize $\cdots$};
\node at (h5) {\scriptsize $i=n$};

\end{tikzpicture}
\caption{$n$ equality gadgets chained together.}
\label{fig:equality-gadget-chain}
\end{figure}

We conclude by noting that the resulting graph $G$ has $\Theta(n)$ vertices and $\Theta(n)$ edges.
\end{proof}

\begin{claim}
\label{claim:disjoint-trav-bound}
Recall that in two-party $\DISJ$, Alice and Bob are given sets $X, Y \subseteq [n]$ each. The goal is to determine whether $X \cap Y = \emptyset$. We claim $D(\TRAV_2) \geq D(\DISJ)$ and so on.
\end{claim}

\begin{proof}
\textbf{Disjointness gadget.}
We construct a gadget that enforces disjointness at a single entry $i \in [n]$. The gadget is shown in Figure \ref{fig:disjoint_gadget}. There, let $x_i$ and $y_1$ be indicator functions whether $i \in X$ and $i \in Y$ respectively. Then:
\begin{enumerate}
    \item The top branch corresponds to when $x_i$ has the element but $y_i$ does not
    \item The middle branch corresponds to when $x_i$ and $y_i$ do not have the element
    \item The top branch corresponds to when $y_i$ has the element but $x_i$ does not
\end{enumerate}
In other words, there is no path left to right if $i \in X \cap Y$.

\begin{figure}[H]
\centering
\begin{tikzpicture}[scale=1.2, every node/.style={circle, fill=black, inner sep=1.5pt}]

\node (L) at (-2,0) {};
\node (T1) at (-0.8,1.5) {};
\node (T2) at (0.8,1.5) {};
\node (M1) at (-0.8,0) {};
\node (M2) at (0.8,0) {};
\node (B1) at (-0.8,-1.5) {};
\node (B2) at (0.8,-1.5) {};
\node (R) at (2,0) {};

\draw (L) -- (T1);
\draw (T1) -- (T2);
\draw (T2) -- (R);

\draw (L) -- (B1);
\draw (B1) -- (B2);
\draw (B2) -- (R);

\draw (L) -- (M1);
\draw (M1) -- (M2);
\draw (M2) -- (R);

\node[draw=none, fill=none, above=6pt] at (T1) {$x_i = 1$};
\node[draw=none, fill=none, above=6pt] at (T2) {$y_i = 0$};

\node[draw=none, fill=none, above=6pt] at (M1) {$x_i = 0$};
\node[draw=none, fill=none, above=6pt] at (M2) {$y_i = 0$};

\node[draw=none, fill=none, below=6pt] at (B1) {$x_i = 0$};
\node[draw=none, fill=none, below=6pt] at (B2) {$y_i = 1$};

\node[draw=none, fill=none, left=6pt] at (L) {$L$};
\node[draw=none, fill=none, right=6pt] at (R) {$R$};

\end{tikzpicture}
\caption{Disjoitness gadget. A path $L$ to $R$ only exists if either $x_i = 0$ or $y_i = 0$.}
\label{fig:disjoint_gadget}
\end{figure}

\textbf{Reduction to our graph problem.}
We construct a graph $G$ by chaining one copy of this gadget for each index $i \in [n]$ between a global start vertex $v$ (left) and end vertex $w$ (right). Construct $G_A$ by giving all vertices at the start and end of every disjointness gadget, and vertices labelled $x_i = 1$ if it is true, else $x_i = 0$. Add all edges between these vertices, and all edges between $x_i = z$ and $y_i = z'$ for $z, z' \in \{0,1\}$. Do similar for $G_B$.

In the resulting graph, there exists a path from $v$ to $w$ if and only if all disjointness gadgets are traversable, which occurs if and only if $X \cap Y = \emptyset$.

\begin{figure}[h]
\centering
\begin{tikzpicture}[
scale=1.5,
transform shape,
hex/.style={
draw,
regular polygon,
regular polygon sides=6,
minimum size=1.4cm,
inner sep=0pt
},
point/.style={circle, fill=black, inner sep=0.5pt}
]

% Hexagons placed manually so edges touch
\node[hex] (h1) at (0,0) {};
\node[hex] (h2) at (1.4,0) {};
\node[hex] (h3) at (2.8,0) {};
\node[hex] (h4) at (4.2,0) {};
\node[hex] (h5) at (5.6,0) {};

% Draw horizontal line through the middle of all hexagons
\draw (-0.7,0) -- (6.3,0);

% Labels moved slightly up
\node[yshift=8pt] at (h1) {\scriptsize $i=1$};
\node[yshift=8pt] at (h2) {\scriptsize $i=2$};
\node[yshift=8pt] at (h3) {\scriptsize $i=3$};
\node[yshift=8pt] at (h4) {\scriptsize $\cdots$};
\node[yshift=8pt] at (h5) {\scriptsize $i=n$};

\node[point, label=left:{\scriptsize $v=L_1$}] at (-0.7,0) {};
\node[point, label=right:{\scriptsize $w=R_n$}] at (6.3,0) {};

\end{tikzpicture}
\caption{Disjointness gadgets chained together.}
\label{fig:disjoint-gadget-chain}
\end{figure}

We conclude by noting that $|V(G)| = \Theta(n)$ and $|E(G)| = \Theta(n)$.
\end{proof}
